\section{Experiments}
\label{sec:experiments}

This section reports the topic-facing evaluation of CR-GNN on protein-related graph datasets. The main goal is not to maximize benchmark count, but to determine whether cross-residual information flow offers a useful advantage over simpler residual baselines under a stable and reproducible protocol.

\subsection{Experimental Settings}

\subsubsection{Model Set}

We evaluate five architectures under the same message-passing backbone:
\begin{enumerate}
    \item \textbf{PlainGNN}: plain stacked propagation
    \item \textbf{NodeResGNN}: single-branch residual reuse
    \item \textbf{NodeCrossGNN}: two-branch node-level cross residual
    \item \textbf{GraphResGNN}: repeated graph-level residual refinement
    \item \textbf{GraphCrossGNN}: graph-level cross residual across sequential block pairs
\end{enumerate}

GraphCondGNN is excluded from the topic-facing main comparison because the present paper aims to sharpen the residual versus cross-residual story rather than to keep every exploratory architecture in the main table.

\subsubsection{Data Protocol}

All topic-facing experiments use the protocol described in Section~\ref{sec:datasets}: stratified 5-fold graph-level evaluation with a stratified validation split inside each training fold. The validation split ratio is 0.1. Model selection is based on validation loss, and the best checkpoint is evaluated on the held-out test fold.

\subsubsection{Optimization}

All runs use Adam with validation-based early stopping, \texttt{ReduceLROnPlateau}, and gradient clipping. The shared defaults are batch size 32, gradient clipping norm 2.0, scheduler factor 0.5, scheduler patience 15, and minimum learning rate $10^{-5}$.

For the plain and residual baselines, we use a common setting:
\begin{itemize}
    \item epochs = 200
    \item patience = 60
    \item learning rate = 0.005
    \item weight decay = $10^{-4}$
    \item dropout = 0.5
    \item hidden dimension = 64
    \item hidden layers = 4
\end{itemize}

The cross-residual models use the stronger settings identified during the tuning stage. In particular, PROTEINS and DD use dataset-specific cross-model hyperparameters, while ENZYMES currently uses the default cross-residual protocol inherited from the tuned V3 setup.

\subsection{Main Topic-Facing Results}

Table~\ref{tab:topic_main_results} reports the main topic-facing benchmark results, and Table~\ref{tab:topic_aux_results} summarizes the associated optimization behavior and parameter counts.

\begin{table}[t]
\centering
\small
\caption{Topic-facing benchmark results on protein-related graph datasets. Results are reported as mean best test accuracy $\pm$ standard deviation over 5 stratified folds.}
\label{tab:topic_main_results}
\begin{tabular}{lccc}
\toprule
Model & PROTEINS & DD & ENZYMES \\
\midrule
PlainGNN & 0.7071 $\pm$ 0.0273 & 0.7216 $\pm$ 0.0163 & 0.2350 $\pm$ 0.0493 \\
NodeResGNN & 0.6926 $\pm$ 0.0457 & \textbf{0.7300 $\pm$ 0.0201} & 0.2450 $\pm$ 0.0704 \\
NodeCrossGNN & 0.6954 $\pm$ 0.0405 & 0.7156 $\pm$ 0.0174 & 0.2850 $\pm$ 0.0528 \\
GraphResGNN & \textbf{0.7196 $\pm$ 0.0420} & 0.7053 $\pm$ 0.0409 & \textbf{0.3067 $\pm$ 0.0779} \\
GraphCrossGNN & 0.6855 $\pm$ 0.0485 & 0.7088 $\pm$ 0.0270 & 0.2917 $\pm$ 0.0732 \\
\bottomrule
\end{tabular}
\end{table}

\begin{table}[t]
\centering
\small
\caption{Optimization and complexity summary on the topic-facing datasets. Parameter counts depend on dataset input dimension, so they are reported per dataset.}
\label{tab:topic_aux_results}
\begin{tabular}{llccc}
\toprule
Dataset & Model & Params & Mean Test Loss & Mean Best Epoch \\
\midrule
PROTEINS & GraphResGNN & 51208 & 0.5769 & 30.0 \\
PROTEINS & PlainGNN & 17026 & 0.6092 & 77.0 \\
PROTEINS & NodeCrossGNN & 33922 & 0.6075 & 82.6 \\
PROTEINS & NodeResGNN & 17026 & 0.6147 & 70.4 \\
PROTEINS & GraphCrossGNN & 68234 & 0.6164 & 57.0 \\
\midrule
DD & NodeResGNN & 22530 & 0.5770 & 25.8 \\
DD & PlainGNN & 22530 & 0.5715 & 41.6 \\
DD & NodeCrossGNN & 36610 & 0.5888 & 27.6 \\
DD & GraphCrossGNN & 90250 & 0.5937 & 23.8 \\
DD & GraphResGNN & 67720 & 0.5865 & 31.4 \\
\midrule
ENZYMES & GraphResGNN & 52248 & 1.7627 & 91.2 \\
ENZYMES & GraphCrossGNN & 69534 & 1.8369 & 73.8 \\
ENZYMES & NodeCrossGNN & 34182 & 1.8442 & 95.2 \\
ENZYMES & NodeResGNN & 17286 & 1.8127 & 24.0 \\
ENZYMES & PlainGNN & 17286 & 1.7904 & 93.4 \\
\bottomrule
\end{tabular}
\end{table}


The results support three practical conclusions.

\textbf{Residual reuse remains the strongest default design.} On PROTEINS, GraphResGNN achieves the best mean accuracy of $0.7196 \pm 0.0420$, improving over PlainGNN ($0.7071 \pm 0.0273$) and both cross-residual variants. On DD, NodeResGNN is best at $0.7300 \pm 0.0201$, again ahead of both NodeCrossGNN and GraphCrossGNN. These two datasets therefore favor residual reuse more clearly than cross-path exchange under the final topic-facing protocol.

\textbf{Cross-residual gains are selective rather than dominant.} ENZYMES is the dataset on which cross-residual variants are most visibly competitive in the final suite: GraphCrossGNN reaches $0.2917 \pm 0.0732$ and NodeCrossGNN reaches $0.2850 \pm 0.0528$, both above PlainGNN ($0.2350 \pm 0.0493$). Even there, however, GraphResGNN remains best at $0.3067 \pm 0.0779$. The practical interpretation is that cross-residual interaction can help on some structure-function classification regimes, but the benefit is not strong enough in the present benchmarks to displace residual baselines as the primary reference.

\textbf{The relevant comparison is structural fit, not architectural complexity alone.} The final benchmark does not identify one universal winner. Instead, it shows that graph-level residual propagation is strongest on PROTEINS and ENZYMES, whereas node-level residual reuse is strongest on DD. This is exactly why the paper emphasizes a controlled family comparison rather than a single headline architecture.

\subsection{Topic-Facing Analysis Assets}

To support interpretation beyond a scalar accuracy number, the current V3 pipeline records:
\begin{itemize}
    \item TensorBoard traces for loss, accuracy, learning rate, gradient norm, and embedding statistics
    \item exported artifacts for graph embeddings, logits, predictions, and confusion-matrix reconstruction
    \item parameter counts and per-fold optimization histories saved in the JSON logs
\end{itemize}

These analysis outputs are useful for checking whether a model improves because it truly preserves useful structure or merely benefits from an unstable training trajectory.

\subsection{Supplementary Benchmarks}

Datasets such as MUTAG, AIDS, and Mutagenicity are still useful as supplementary robustness checks, but they are not used to define the main biological claim of the paper. MSRC\_9 is omitted from the main text because it is outside the biological scope of the target venue. This separation is intentional: the topic-facing narrative should be driven by protein-related evidence rather than by generic graph benchmark breadth.

\subsection{Current Limitations}

The experiments are still structure-centric rather than fully omics-integrative. As a result, the present paper supports a conservative claim: CR-GNN is an effective biomolecular graph representation architecture on protein-related benchmarks. It does not yet justify a stronger claim about integrative omics prediction or direct discovery of novel genes, domains, or motifs. That stronger claim would require either node-level biological prediction tasks or graph datasets with richer heterogeneous biological signals.
